

\chapter{Overview: Fitting diffraction patterns}

There are four key aspects for obtaining a good fit to a powder diffraction pattern in a Rietveld refinement: 
\begin{itemize}
     
\item Fitting the background, which means reproducing the slowly-varying parts of the diffraction pattern, that arise from instrumental artifacts as well as non-Bragg scattering from the sample (Chapter \ref{Chap:background}),

\item Fitting the peak positions, which are determined by lattice parameters and instrumental calibration/corrections must be in the correct places
(Chapter \ref{Chap:peakpositions}).


\item Reproducing the peak intensities, as determined from the crystal structure model as well as corrections such as preferred orientation and extinction 
(Chapter \ref{Chap:intensities}).


\item Matching the shapes of the peaks, as this must reproduce the broadening effects of both the diffraction instrument and effects from the sample (Chapter \ref{Chap:profile}). 

\end{itemize}

Each of these aspects of fitting will be discussed in the following chapters. 
